% ---------------------------------------------------------------------------
% content/build.tex -- how compile_thesis.sh works and why.
% ---------------------------------------------------------------------------

\linkedsection{script}{What the Script Does}

\texttt{compile\_thesis.sh} needs TeX Live (or MacTeX) with \texttt{latexmk}, \texttt{pdflatex}, \texttt{xelatex} and \texttt{biber}, and Python~3 with the \texttt{pikepdf} package.
It runs the three steps of \Cref{fig:pipeline}, all inside the git-ignored folder \texttt{build/}.

\begin{figure}[htb]
\hypertarget{target:fig:pipeline}{}%
    \centering
    \resizebox{\textwidth}{!}{%
    \begin{tikzpicture}[
        box/.style={draw,rounded corners,minimum height=9mm,minimum width=27mm,align=center,font=\small},
        arr/.style={-{Stealth[length=2.5mm]},thick}]
        \node[box] (en) {\texttt{main.tex}\\pdfLaTeX + biber};
        \node[box,below=9mm of en] (he) {\texttt{hebrew\_pages.tex}\\XeLaTeX + biber};
        \node[box,right=22mm of en] (enpdf) {\texttt{build/main.pdf}};
        \node[box,right=22mm of he] (hepdf) {\texttt{build/hebrew\_pages.pdf}};
        \node[box,right=22mm of $(enpdf)!0.5!(hepdf)$] (merge) {\texttt{merge\_thesis.py}};
        \node[box,right=12mm of merge] (out) {\texttt{build/thesis.pdf}};
        \draw[arr] (en) -- (enpdf);
        \draw[arr] (he) -- (hepdf);
        \draw[arr] (enpdf) -- (merge);
        \draw[arr] (hepdf) -- (merge);
        \draw[arr] (merge) -- (out);
        \draw[arr,dashed] (en.south) to[bend right=25] node[left,font=\scriptsize,align=right] {page numbers,\\\texttt{LastPage}} (he.north);
        \draw[arr,dashed] (he.north east) to[bend right=25] node[right,font=\scriptsize,align=left] {Hebrew\\page numbers} (en.south east);
    \end{tikzpicture}}
    \caption{The build. Solid arrows are the three steps of the script; dashed arrows are the labels each document reads from the other's auxiliary file, which is why the first two steps are repeated until both settle.}
    \label{fig:pipeline}
\end{figure}

\begin{enumerate}
    \item \texttt{latexmk -pdf main.tex} produces \texttt{build/main.pdf}. \texttt{latexmk} runs biber and repeats pdfLaTeX until the references are stable.
    \item \texttt{latexmk -xelatex hebrew\_pages.tex} produces \texttt{build/hebrew\_pages.pdf}.
    \item \texttt{merge\_thesis.py} appends the Hebrew pages to the English PDF and rewires the cross-document links (\Cref{sec:merge}). The result is \texttt{build/thesis.pdf}.
\end{enumerate}

\texttt{./compile\_thesis.sh clean} deletes \texttt{build/}. A clean build takes a few minutes; later builds only redo what changed.

\linkedsection{settle}{Why the First Two Steps Repeat}

Each document reads labels from the other.
The Hebrew document continues the page numbering after the English \texttt{LastPage} label and prints the English chapter numbers and pages in its table of contents.
The English document, in turn, needs the page numbers of the Hebrew summary and abstract for its own table of contents.
Both documents load the \texttt{xr-hyper} package, which imports the labels of another document's \texttt{.aux} file under a prefix: \texttt{en-} for English labels read by the Hebrew document, \texttt{he-} for Hebrew labels read by the English one.

Because each side depends on the other, the script compiles English, then Hebrew, then checks whether either \texttt{.aux} file changed, and repeats the pair until neither does.
On a clean build the first English pass prints page~0 for the Hebrew entries, since no Hebrew \texttt{.aux} exists yet; the second pass fixes it.
From the second pass on, the script passes \texttt{-g} to \texttt{latexmk}, because \texttt{latexmk} does not know that the other document's \texttt{.aux} file is a dependency and would otherwise report that there is nothing to do.

\linkedsection{troubleshooting}{Fonts, Overleaf and Troubleshooting}

\paragraph{Fonts.} The Hebrew pages use \emph{Linux Libertine O} for Latin text, \emph{David CLM} for Hebrew and \emph{Latin Modern Roman} for Latin fragments inside Hebrew paragraphs. All three ship with TeX Live (packages \texttt{libertine}, \texttt{culmus}, \texttt{lm}), but XeTeX on macOS only finds a font by name if macOS itself knows it. If the build stops with \texttt{The font "Linux Libertine O" cannot be found}, run \texttt{./compile\_thesis.sh install-fonts} once; it copies the font files into \texttt{\textasciitilde/Library/Fonts}. On Linux and on Overleaf this is not needed.

\paragraph{Overleaf.} Overleaf cannot run the merge script. Compile \texttt{main.tex} with pdfLaTeX and \texttt{hebrew\_pages.tex} with XeLaTeX (two projects, or one project with the main file switched), download both PDFs and merge locally with \texttt{python3 merge\_thesis.py english.pdf hebrew.pdf out.pdf}.

\paragraph{Symptoms and causes.}
\begin{itemize}
    \item \emph{Hebrew Summary~0 in the English TOC}: the English part was not rebuilt after the Hebrew part. Run the script again rather than \texttt{latexmk} by hand.
    \item \emph{Merge warning about names that exist in neither PDF}: a Hebrew \verb|\hyperlink| has no matching anchor in the English source. See \Cref{sec:merge}.
    \item \emph{Nonsense errors such as \texttt{Missing \textbackslash begin\{document\}} in \texttt{build/main.aux}}: a build was interrupted or two ran at once and an \texttt{.aux} file is truncated. Run \texttt{./compile\_thesis.sh clean} and build again.
    \item \emph{biber: extracting arm64 binary with lipo failed} (macOS, after a system update): the universal \texttt{biber} shipped with MacTeX unpacks itself with \texttt{lipo}, which newer command line tools reject. Extract the slice once by hand, \texttt{lipo -thin arm64 \$(which biber) -output \textasciitilde/bin/biber}, and put \texttt{\textasciitilde/bin} first on your \texttt{PATH}.
    \item \emph{Package bidi error about xcolor}: in \texttt{hebrew\_pages.tex}, any package that loads \texttt{xcolor} must come before \texttt{bidi} and \texttt{polyglossia}. The template loads \texttt{xcolor} first for this reason.
\end{itemize}
